% !TeX root = ../main.tex
%%%%%%%%%%%%%%%%%%%%%%%%%%   5   %%%%%%%%%%%%%%%%%%%%%%%%%%%%%%
\section{如期完成全部论文工作的可能性}

从已有基础看，目前已完成非线性滤波与最优输运的基础理论整理、代表性算法实现及初步对照实验，形成了可供后续研究使用的推导资料、程序与评价基线。这些工作为进一步学习神经网络滤波方法、选择研究问题和验证算法提供了基础，有助于减少重复准备工作。

从后续安排看，9--10 月以文献学习和数学原理总结为重点，11--12 月围绕选定的仿真环境开展算法设计与实验验证，前后阶段具有明确的衔接关系。研究过程中同步整理文献综述、方法推导与实验分析，逐步形成论文内容，避免将资料整理和论文撰写全部集中于研究末期。

主要风险仍在于理论理解、算法调试和实验分析所需时间存在不确定性。为此，拟将研究范围限定于能够充分分析和验证的具体问题，依据每日记录与阶段总结评估实际进展，优先完成支撑论文核心论证的工作，并为结果复核和论文修改预留时间。若实验效果未达预期，应及时分析原因与适用边界，调整技术方案，避免在缺乏证据支持的方向上持续投入。

综合而言，现有工作已具备继续推进的基础；在合理控制研究范围、落实阶段目标并根据实际进展及时调整安排的前提下，如期完成全部论文工作具有可行性。
